\documentclass[
  spanish,
  aspectratio=169,
  xcolor={dvipsnames,table}
]{beamer}

\geometry{paperwidth=19.2cm,paperheight=10.8cm}
\usepackage[utf8]{inputenc}
\usepackage[T1]{fontenc}
\usepackage[spanish,es-tabla]{babel}
\usepackage[scaled=.96]{helvet}
\renewcommand{\familydefault}{\sfdefault}
\usepackage{booktabs}
\usepackage{graphicx}
\usepackage{ragged2e}
\usepackage{tikz}
\usepackage{hyperref}
\usetikzlibrary{calc,positioning,fit,shadows.blur}

\newcommand{\studentname}{Martin Jonathan de la Cruz}
\newcommand{\studentshort}{M. J. de la Cruz}
\newcommand{\studentid}{ES2611202040}
\newcommand{\universityname}{Universidad Abierta y a Distancia de México}
\newcommand{\facultyname}{Licenciatura en Derecho}
\newcommand{\coursename}{Garantías Constitucionales}
\newcommand{\coursecode}{LDE-S2B1}
\newcommand{\activitytitle}{El Estado a través de la historia}
\newcommand{\activitysubtitle}{Foro: poder, derecho y garantías}
\newcommand{\documentsubject}{Actividad 2 - Garantías Constitucionales}
\newcommand{\activityweek}{Semana 2}
\newcommand{\teachingfigure}{Aarón López Pérez}
\newcommand{\deliverydate}{\today}
\newcommand{\locationname}{Roma Norte, Ciudad de México}
\newcommand{\departmentlogo}{img/departamentos/UnADM.pdf}

\definecolor{unadmGreenDark}{HTML}{174A3A}
\definecolor{unadmGreen}{HTML}{5F8F3A}
\definecolor{unadmGold}{HTML}{B88A2A}
\definecolor{unadmPaper}{HTML}{F6F7F2}
\definecolor{unadmInk}{HTML}{1F2A24}

\mode<presentation>{
  \usetheme{default}
  \usefonttheme{professionalfonts}
  \setbeamertemplate{navigation symbols}{}
  \setbeamertemplate{blocks}[rounded][shadow=false]
}
\setbeamersize{text margin left=0.60cm,text margin right=0.60cm}
\setbeamercolor{background canvas}{bg=white}
\setbeamercolor{normal text}{fg=unadmInk,bg=white}
\setbeamercolor{structure}{fg=unadmGreenDark}
\setbeamercolor{frametitle}{fg=white,bg=unadmGreenDark}
\setbeamercolor{block title}{fg=white,bg=unadmGreen}
\setbeamercolor{block body}{fg=unadmInk,bg=unadmPaper}
\setbeamerfont{title}{size=\LARGE,series=\bfseries}
\setbeamerfont{frametitle}{size=\large,series=\bfseries}
\setbeamertemplate{itemize item}{\textcolor{unadmGreen}{\large$\blacktriangleright$}}

\setbeamertemplate{frametitle}{
  \nointerlineskip
  \begin{beamercolorbox}[wd=\textwidth,ht=1.02cm,dp=0.22cm,leftskip=0.35cm,rightskip=0.25cm]{frametitle}
    \insertframetitle\hfill\includegraphics[height=0.60cm]{\departmentlogo}
  \end{beamercolorbox}
  {\color{unadmGold}\rule{\textwidth}{1.3pt}}
}

\setbeamertemplate{footline}{
  \leavevmode
  \hbox{
    \begin{beamercolorbox}[wd=.34\paperwidth,ht=2.7ex,dp=1ex,leftskip=1em]{author in head/foot}\color{white}\studentshort\end{beamercolorbox}
    \begin{beamercolorbox}[wd=.40\paperwidth,ht=2.7ex,dp=1ex,center]{title in head/foot}\color{white}\coursename\end{beamercolorbox}
    \begin{beamercolorbox}[wd=.26\paperwidth,ht=2.7ex,dp=1ex,rightskip=1em plus 1fil]{date in head/foot}\color{white}\coursecode\hfill\insertframenumber/\inserttotalframenumber\end{beamercolorbox}
  }
}
\setbeamercolor{author in head/foot}{bg=unadmGreenDark}
\setbeamercolor{title in head/foot}{bg=unadmGreen}
\setbeamercolor{date in head/foot}{bg=unadmGreenDark}

\title[\coursecode]{\activitytitle}
\subtitle{\activitysubtitle}
\author[\studentshort]{\studentname\\\footnotesize Matrícula: \studentid}
\institute[UnADM]{\universityname\\\facultyname}
\date[\activityweek]{\locationname\\\activityweek\\\deliverydate}

\begin{document}

\begin{frame}[plain]
  \begin{tikzpicture}[remember picture,overlay]
    \fill[unadmGreenDark] (current page.south west) rectangle ([xshift=0.58\paperwidth]current page.north west);
    \fill[unadmGreen] ([xshift=0.58\paperwidth]current page.south west) rectangle (current page.north east);
    \node[anchor=center,opacity=0.10] at (current page.center) {\includegraphics[width=0.72\paperwidth]{\departmentlogo}};
    \node[anchor=north west] at ([xshift=0.58cm,yshift=-0.46cm]current page.north west) {\includegraphics[width=2.55cm]{\departmentlogo}};
  \end{tikzpicture}
  \vspace*{1.62cm}
  {\color{white}\fontsize{24}{29}\selectfont\bfseries \activitytitle\par}
  \vspace{0.28cm}
  {\color{white!86}\large \activitysubtitle}\\[0.16cm]
  {\color{white!82}\normalsize \documentsubject}\\[0.35cm]
  {\color{unadmGold}\rule{0.58\linewidth}{1.3pt}}\\[0.35cm]
  {\color{white}\studentname}\\
  {\color{white!82}\footnotesize \facultyname}
\end{frame}

\begin{frame}{Objetivo del foro}
  \begin{block}{Propósito analítico}
    Analizar la evolución histórica del Estado para explicar cómo la limitación jurídica del poder favoreció el reconocimiento y la protección de los derechos humanos.
  \end{block}
  \begin{itemize}
    \item Recorrer las etapas del Estado desde la Antigüedad hasta el modelo constitucional.
    \item Relacionar cada etapa con el avance o retroceso de los derechos.
    \item Sostener una postura personal fundamentada en fuentes académicas.
  \end{itemize}
\end{frame}

\begin{frame}{El Estado en la Antigüedad}
  \begin{columns}[T]
    \begin{column}{0.48\linewidth}
      \begin{block}{Ciudad-Estado griega}
        Ciudadanía restringida; esclavitud naturalizada. Los derechos dependían de la pertenencia a la polis.
      \end{block}
    \end{column}
    \begin{column}{0.48\linewidth}
      \begin{block}{Imperio romano}
        Ius civile para ciudadanos; ius gentium como puente hacia un derecho más universal.
      \end{block}
    \end{column}
  \end{columns}
  \vspace{0.25cm}
  \begin{block}{Rasgo común}
    El poder no reconocía límites jurídicos formales: la protección dependía del estatus social.
  \end{block}
\end{frame}

\begin{frame}{Estado medieval y absolutismo}
  \begin{itemize}
    \item \textbf{Feudalismo}: el poder se fragmenta en señoríos; el campesino carece de derechos efectivos.
    \item \textbf{Monarquía absoluta}: el rey concentra toda autoridad, sin contrapesos institucionales.
    \item \textbf{Iglesia}: actúa como freno moral al poder secular, pero no garantía jurídica.
  \end{itemize}
  \begin{block}{Magna Carta 1215}
    Primera limitación escrita del poder real: protege derechos de los barones, semilla del constitucionalismo.
  \end{block}
\end{frame}

\begin{frame}{Iusnaturalismo y contractualismo}
  \begin{columns}[T]
    \begin{column}{0.48\linewidth}
      \begin{block}{Iusnaturalismo}
        Derechos anteriores al Estado, derivados de la razón o de la dignidad humana (Locke, Rousseau).
      \end{block}
    \end{column}
    \begin{column}{0.48\linewidth}
      \begin{block}{Contrato social}
        El Estado surge para proteger derechos previos; si los viola, pierde legitimidad.
      \end{block}
    \end{column}
  \end{columns}
  \vspace{0.25cm}
  Esta transformación conceptual hizo posible las declaraciones de derechos del siglo XVIII.
\end{frame}

\begin{frame}{Revoluciones y declaraciones de derechos}
  \begin{itemize}
    \item \textbf{Revolución Gloriosa (1688)}: Bill of Rights limita al monarca y afirma prerrogativas parlamentarias.
    \item \textbf{Declaración de Independencia (1776)}: vida, libertad y búsqueda de la felicidad como derechos inalienables.
    \item \textbf{Declaración de Derechos del Hombre (1789)}: libertad, igualdad, soberanía nacional, separación de poderes.
  \end{itemize}
  \begin{block}{Tesis central}
    Cada revolución redujo el margen del poder arbitrario y amplió el reconocimiento de derechos.
  \end{block}
\end{frame}

\begin{frame}{Estado liberal y Estado social}
  \begin{columns}[T]
    \begin{column}{0.48\linewidth}
      \begin{block}{Estado liberal (siglo XIX)}
        Derechos civiles y políticos. El Estado no interviene; garantiza la esfera privada.
      \end{block}
    \end{column}
    \begin{column}{0.48\linewidth}
      \begin{block}{Estado social (siglo XX)}
        Derechos económicos, sociales y culturales. El Estado actúa para reducir desigualdad.
      \end{block}
    \end{column}
  \end{columns}
  \vspace{0.25cm}
  La Constitución mexicana de 1917 es pionera en reconocer derechos sociales con rango constitucional.
\end{frame}

\begin{frame}{Estado constitucional y garantías}
  \begin{block}{Postura personal}
    El Estado constitucional representa la forma más elaborada de limitación del poder: los derechos se protegen mediante garantías exigibles y mecanismos de control como el juicio de amparo.
  \end{block}
  \begin{itemize}
    \item Constitución como norma suprema vinculante.
    \item Garantías como instrumentos que hacen efectivos los derechos.
    \item Control de constitucionalidad como freno al poder arbitrario.
  \end{itemize}
\end{frame}

\begin{frame}{Fuentes consultadas}
  \begin{itemize}
    \item Rabinovich-Berkman, R. --- \textit{¿Cómo se hicieron los Derechos Humanos?}
    \item Marcano Salazar, L. M. --- \textit{Derechos humanos: teorías y doctrina}
    \item Álvarez Icaza Longoria, E. --- \textit{Los derechos humanos en México}
    \item Nota de estudio: \textit{El Estado a través de la historia} (corpus local)
  \end{itemize}
  \begin{block}{Perspectiva de cierre}
    Cada etapa histórica confirmó que solo al limitar el poder con normas exigibles emergieron derechos humanos reales, no meramente declarados.
  \end{block}
\end{frame}

\begin{frame}{Pregunta para el intercambio}
  \begin{block}{Invitación al debate}
    ¿Puede el Estado constitucional contemporáneo retroceder hacia formas de poder menos limitadas, o los mecanismos de garantía y el derecho internacional hacen ese retroceso estructuralmente imposible?
  \end{block}
  \vspace{0.35cm}
  \begin{itemize}
    \item Reflexiona sobre la solidez de las garantías actuales.
    \item Contrasta con los retrocesos históricos ya estudiados.
    \item ¿Qué diferencia hay entre un derecho reconocido y uno garantizado?
  \end{itemize}
\end{frame}

\begin{frame}{Cierre}
  \begin{block}{Aprendizaje principal}
    La evolución del Estado no es lineal ni automática: cada conquista jurídica fue resultado de luchas sociales y transformaciones institucionales que redujeron el poder arbitrario.
  \end{block}
  \begin{itemize}
    \item Las garantías constitucionales son el resultado histórico de siglos de limitación del poder.
    \item Entender el Estado es condición para defender los derechos que lo vinculan.
    \item El foro permite contrastar estas perspectivas con otras visiones del grupo.
  \end{itemize}
\end{frame}

\end{document}
